\documentclass{article}
\usepackage[T1]{fontenc}
\usepackage{lmodern}
\usepackage{graphicx}
\usepackage{amsmath,amssymb}
\usepackage[a4paper,margin=2cm]{geometry}
\usepackage[absolute,overlay]{textpos}
\setlength{\TPHorizModule}{1mm}
\setlength{\TPVertModule}{1mm}
\usepackage[indonesian]{babel}
\usepackage{hyperref}
\setlength{\emergencystretch}{2em}
% unit: Halaman 10

\begin{document}

% === Halaman 10 (llm) ===

\section*{Pembangkitan Kunci Internal (\textit{key expansion})}

\begin{itemize}
    \item Kunci internal = kunci setiap putaran = \textit{subkey (round key)}
    \item Ada 16 putaran, jadi ada 16 kunci internal: $K_1, K_2, \dots, K_{16}$
    \item Kunci internal dibangkitkan dari kunci eksternal (64 bit) yang diberikan oleh pengguna. Di dalam dokumen resmi DES pembangkitan kunci internal disebut \textit{key schedule}.
    \item Panjang kunci internal adalah 48 bit
    \item Gambar 3 memperlihatkan proses pembangkitan kunci internal.
\end{itemize}

\end{document}
